\section*{ID10142: Phase 1: Proof of Concept (Current Technology): μ\_r ≈ 5 × 10\textasciicircum{}4}
\paragraph{Metadata.}
PiID: 6655566156787 | RTEID: RTE:P39383.6383:T2.0603 | UUID: 49aa23a0-c15d-54e8-8b62-74f4c0a9280a

\paragraph{Canonical Statement.}
\textbackslash{}textbf\{Components:\} \textbackslash{}begin\{itemize\} \textbackslash{}item Gyroscope: Ceramic hybrid bearings, \$\textbackslash{}omega = 5 \textbackslash{}times 10\textasciicircum{}4\$ RPM \textbackslash{}item Dipoles: NdFeB permanent magnets, \$B\_r = 1.4\$ T \textbackslash{}item Boundary: Mu-metal sphere, \$\textbackslash{}mu\_r \textbackslash{}approx 5 \textbackslash{}times 10\textasciicircum{}4\$ \textbackslash{}item Sensors: Commercial SQUID, \$10\textasciicircum{}\{-14\}\$ T sensitivity \textbackslash{}end\{itemize\}

\paragraph{Canonical Equation.}
\[
\mu_r \approx 5 \times 10^4
\]

\paragraph{Dictionary.}
Phase 1: Proof of Concept (Current Technology): μ\_r ≈ 5 × 10\textasciicircum{}4 — μ\_r ≈ 5 × 10\textasciicircum{}4

\paragraph{Encyclopedia.}
Phase 1: Proof of Concept (Current Technology): μ\_r ≈ 5 × 10\textasciicircum{}4 is a scaling / order-of-magnitude relation, written μ\_r ≈ 5 × 10\textasciicircum{}4. It expresses a scaling or proportionality, capturing how the quantity grows with its parameters. It is built from r. Within the Trawin Zero Point registry it serves as one element of the framework's mathematical narrative.
